\section{Laboratory Details}
\label{sec:appendix}

\subsection{Generated-Traffic Lab}

\begin{table}[htbp]
  \centering
  \caption{Components of the generated-traffic lab.}
  \label{tab:lab-components}
  \footnotesize
  \begin{tabularx}{\linewidth}{@{}llY@{}}
    \toprule
    \textbf{Component} & \textbf{Address} & \textbf{Role} \\
    \midrule
    Target & 172.28.0.10 & nginx web server in front of a Python application \\
    Attacker & 172.28.0.20 & Generates the HTTP flood and the SYN flood \\
    Benign probe & 172.28.0.30 & Requests the page continuously to measure availability and latency \\
    Sniffer & Shares the target's network namespace & Captures all traffic on the target's interface \\
    \bottomrule
  \end{tabularx}
\end{table}

nginx served \path{/heartbeat} as a static file and forwarded \path{/} and \path{/index.html} to a local Python application. Each page request read and sorted a 5,000-row catalogue before responding, giving a repeatable workload. A local health check ran every two seconds; it appeared in the nginx logs but not in the capture.

The target was limited to half a CPU core and 64 MiB of memory, inside a virtual machine with two vCPUs, about 3.7 GiB of memory and no swap. The other containers had no resource limits. SYN cookies were disabled on the target.

\begin{table}[htbp]
  \centering
  \caption{Tools used in the generated-traffic lab.}
  \label{tab:lab-tools}
  \begin{tabularx}{\linewidth}{@{}lY@{}}
    \toprule
    \textbf{Tool} & \textbf{Purpose} \\
    \midrule
    nginx and Python & Web server and application workload on the target \\
    hey & HTTP flood generation \\
    ZMap & SYN flood generation \\
    curl & Benign probe requests and latency measurement \\
    tcpdump & Packet capture on the target's interface \\
    \bottomrule
  \end{tabularx}
\end{table}

\subsection{Real-Traffic Lab}

\begin{table}[htbp]
  \centering
  \caption{Real-traffic dataset.}
  \label{tab:real-dataset}
  \footnotesize
  \begin{tabularx}{\linewidth}{@{}lY@{}}
    \toprule
    \textbf{Item} & \textbf{Value} \\
    \midrule
    File & \path{2025-01-22-traffic-analysis-exercise.pcap} \\
    Size & 26,759,675 bytes; 39,427 packets; 19:44:56 to 20:38:18 UTC \\
    SHA-256 & \texttt{\footnotesize e59db1c07c6fdefafa0abdbca03248c341cdc36c09c34753204d3162802a3586} \\
    Scenario & A user in an Active Directory network searched for Google Authenticator, downloaded a fake installer and the host then contacted attacker infrastructure \\
    Network & 10.1.17.0/24, domain bluemoontuesday.com, domain controller 10.1.17.2 \\
    \bottomrule
  \end{tabularx}
\end{table}

The Docker Compose stack ran on a single host with four CPUs and 6 GB of memory.

\begin{table}[htbp]
  \centering
  \caption{Components of the real-traffic lab.}
  \label{tab:real-components}
  \footnotesize
  \begin{tabularx}{\linewidth}{@{}lY@{}}
    \toprule
    \textbf{Component} & \textbf{Role} \\
    \midrule
    Zeek & Converts packets into connection, DNS, HTTP, TLS, certificate and file logs \\
    Suricata & Matches packets against ET Open and the group's local rules \\
    Filebeat & Ships Zeek and Suricata logs into Elasticsearch in a common schema \\
    Elasticsearch and Kibana & Store and search events; host the security dashboard \\
    Flow-statistics script & Scores beaconing and HTTP request rhythm from the Zeek logs \\
    tshark, Wireshark and Python & Packet inspection and the independent flow table \\
    \bottomrule
  \end{tabularx}
\end{table}
